\section{Detailed scoring and intervention construction}
\label{app:detailed-method}

Our method compares how different vision--language models (VLMs) use localized visual evidence without assuming that their internal coordinates or sparse features are homologous. For each internal-object comparison, we hold the image--question input, evidence corruption, target answer, and answer readout fixed, then vary the object intervened on. This common input-side reference lets us evaluate dense states, fitted sparse reconstructions, native channel groups, and model-specific residual-state subspaces shared across examples without equating their coordinates or claiming that every object is tested on one identical dataset.

The method tests two hypotheses. The \textbf{regional-specificity hypothesis} predicts that corrupting a human-annotated answer-bearing region lowers the target-answer score more than matched non-evidence corruption, and that restoring aligned clean internal states recovers more score under evidence corruption than under matched control corruption. The \textbf{path-preservation hypothesis} asks whether a qualifying internal object can recover the score change caused by masking the evidence and remove a corresponding change when its corrupted value replaces the clean value. Compact sparse, native-channel, and shared-subspace candidates are tested before the complete dense state. Structural differences alone do not satisfy either requirement. This section first defines answer scores and matched causal contrasts, then the internal objects and intervention operations used to test the two hypotheses. Section~\ref{sec:experimental-setup} subsequently fixes the models, evaluation sets, selection roles, and statistical criteria.

\subsection{Evaluation conditions and answer support}

Our comparisons evaluate the same input $x$ and target answer $a$ under different complete forward-run configurations. Let $\mathcal{C}$ be the set of these evaluation conditions. Each $c\in\mathcal{C}$ specifies whether the image is unchanged or corrupted and whether selected internal states are unchanged or intervened on; it is an evaluation configuration rather than an additional text input. To score the same answer in every condition, we provide the fixed target prefix $a_{<t}$ before predicting each token rather than feeding back a condition-specific generated prefix. This teacher-forced forward pass produces condition-specific states $\vec{h}_{i}^{\ell}(x,c;\theta)$ and next-token logit vectors $\vec{z}_{t}(x,c;\theta)$. The target answer and its tokenization remain fixed across paired conditions, so the prefix is omitted from these condition-specific argument lists.

The probability assigned to target token $a_t$ under condition $c$ is
\begin{equation}
    p_t^{c}
    =\frac{\exp\!\left([\vec{z}_{t}(x,c;\theta)]_{a_t}\right)}
           {\sum_{v\in\mathcal{V}}\exp\!\left([\vec{z}_{t}(x,c;\theta)]_{v}\right)},
    \label{eq:target-token-probability}
\end{equation}
where $v$ indexes the tokens in $\mathcal{V}$. We use the answer-support set $\mathcal{S}=\{S_{\mathrm{logit}},S_{\mathrm{seq}},S_{\mathrm{rank}}\}$. The first-token logit score is
\begin{equation}
    S_{\mathrm{logit}}(x,a;c)
    =[\vec{z}_{1}(x,c;\theta)]_{a_1},
    \label{eq:logit-score}
\end{equation}
the complete-answer sequence score is
\begin{equation}
    S_{\mathrm{seq}}(x,a;c)
    =\sum_{t=1}^{T}\log p_t^{c},
    \label{eq:app-sequence-score}
\end{equation}
\begin{samepage}
The first-token rank score is
\begin{equation}
    S_{\mathrm{rank}}(x,a;c)
    =-\mathtt{Rank}\!\left(a_1,\vec{z}_{1}(x,c;\theta),\mathcal{V}\right).
    \label{eq:rank-score}
\end{equation}
The operation $\mathtt{Rank}$ returns the one-based descending rank of $a_1$ among the tokens in $\mathcal{V}$. The negative sign gives all three scores the same direction: larger values mean stronger support for the target answer. For a matched wrong answer $\bar{a}$ to the same input, the correct-versus-wrong margin is
\begin{equation}
    M_S(x,a,\bar{a};c)
    =S(x,a;c)-S(x,\bar{a};c),
    \qquad S\in\mathcal{S}.
    \label{eq:answer-margin}
\end{equation}
A positive margin means that the model assigns greater support to the correct target $a$ than to the matched wrong answer $\bar{a}$.
\end{samepage}

\subsection{Matched causal contrasts}

Let $c_0\in\mathcal{C}$ be the clean condition, $c_{\mathrm{evid}}\in\mathcal{C}$ an evidence-region corruption, and $c_{\mathrm{ctrl}}\in\mathcal{C}$ a matched control corruption. For any answer-support function $S\in\mathcal{S}$ defined above, the damage caused by condition $c$ is
\begin{equation}
    D_S(x,a;c)=S(x,a;c_0)-S(x,a;c).
    \label{eq:app-corruption-damage}
\end{equation}
Larger $D_S$ means that the corruption removes more support for the fixed target answer $a$. The evidence-specific regional contrast is
\begin{equation}
    E_S(x,a;c_{\mathrm{evid}},c_{\mathrm{ctrl}})
    =D_S(x,a;c_{\mathrm{evid}})-D_S(x,a;c_{\mathrm{ctrl}}).
    \label{eq:app-evidence-control-contrast}
\end{equation}
Thus, $E_S>0$ means that corrupting the annotated evidence is more consequential than corrupting its matched control. The clean-to-evidence-mask damage $D_S(x,a;c_{\mathrm{evid}})$ is a different quantity: $E_S$ subtracts the damage from the control mask.

We next connect this input-side deficit to an internal model state. Let $\mathcal{O}$ be the set of intervention objects, where an object $o\in\mathcal{O}$ can be a complete hidden state, a projected hidden-state difference, a selected sparse reconstruction, a reconstruction residual, or a native-channel group. The condition $\mathtt{Restore}(o,c)$ replaces object $o$ in corrupted run $c$ with its clean value. Its restoration benefit is
\begin{equation}
    B_S(x,a;o,c)
    =S\!\left(x,a;\mathtt{Restore}(o,c)\right)-S(x,a;c).
    \label{eq:app-restoration-benefit}
\end{equation}
Restoration tests whether the selected object is sufficient to recover the target-answer score lost under corruption. The complementary condition $\mathtt{Delete}(o,c)$ is shorthand for corrupted-state replacement: it writes the object's value from corrupted run $c$ into the clean run $c_0$, rather than zeroing the object. The resulting deletion damage is
\begin{equation}
    A_S(x,a;o,c)
    =S(x,a;c_0)-S\!\left(x,a;\mathtt{Delete}(o,c)\right).
    \label{eq:app-deletion-damage}
\end{equation}
Deletion damage measures whether replacing the selected object's clean value with its corrupted value removes the clean-versus-corrupted target-score difference. This is a test of necessity under the specified corrupted-state replacement, not under every possible ablation. The comparisons specific to each result evaluate $B_S$ and $A_S$ against equal-norm random directions, reversed clean-state differences, wrong insertion positions, or full-channel references. Together, $E_S$ and these internal controls separate evidence-specific input damage from score changes produced by restoring or replacing one internal object. Restoration alone does not test whether the clean run needs that object's clean value; corrupted-state replacement alone does not test whether restoring it recovers the masked deficit. An object passes the bidirectional interface criterion only if both directions meet their respective gates.

\subsection{Dense and sparse representation changes}

Applying Equation~\ref{eq:sparse-encoding} to a condition-specific dense state defines $f_{i,j}^{\ell}(x,c;\theta,\phi)$. Let $\vec{u}_{i}^{m}(x,c;\theta)\in\mathbb{R}^{d}$ be the complete feed-forward output written at position $i$ in target layer $m$. For a visually corrupted condition $c$, the prefix $\Delta$ denotes a clean-minus-corrupted difference at the same internal object and position:
\begin{align}
    \Delta\vec{h}_{i}^{\ell}(x,c;\theta)
    &=\vec{h}_{i}^{\ell}(x,c_0;\theta)-\vec{h}_{i}^{\ell}(x,c;\theta),
    \label{eq:hidden-difference}\\
    \Delta\vec{u}_{i}^{m}(x,c;\theta)
    &=\vec{u}_{i}^{m}(x,c_0;\theta)-\vec{u}_{i}^{m}(x,c;\theta),
    \label{eq:feedforward-difference}
\end{align}
The first quantity is the residual-state change used for hidden restoration. The second is the complete feed-forward-output change used for the same-object sparse comparison.

Let $\mathcal{J}^{\ell}=\{1,\ldots,n_{\mathrm{feat}}^{\ell}\}$ be the complete feature-index set in read layer $\ell$. For a fixed budget $K$, let $\mathcal{K}_{i}^{\ell}(x,c;K)\subseteq\mathcal{J}^{\ell}$ be the indices of the $K$ largest nonnegative feature activations at position $i$ under condition $c$. The clean and corrupted runs select their top-$K$ indices separately. For feature $j$ read in layer $\ell$, let $\vec{w}_{j}^{\ell\rightarrow m}(\phi)\in\mathbb{R}^{d}$ be its decoder direction into target layer $m\in\{\ell,\ldots,L-1\}$. A PLT uses the within-layer case $m=\ell$; a CLT can write to a later target layer $m>\ell$. The selected-feature reconstruction of the clean--corrupted change at position $i$ is
\begin{equation}
\begin{aligned}
    \widehat{\Delta\vec{u}}_{K,i}^{\ell\rightarrow m}(x,c;\theta,\phi)
    ={}&\sum_{j\in\mathcal{K}_{i}^{\ell}(x,c_0;K)}
      f_{i,j}^{\ell}(x,c_0;\theta,\phi)\,\vec{w}_{j}^{\ell\rightarrow m}(\phi)\\
    &-\sum_{j\in\mathcal{K}_{i}^{\ell}(x,c;K)}
      f_{i,j}^{\ell}(x,c;\theta,\phi)\,\vec{w}_{j}^{\ell\rightarrow m}(\phi).
\end{aligned}
    \label{eq:sparse-reconstruction}
\end{equation}
The hat marks a reconstruction rather than a native model state. Any condition-independent affine decoder offset cancels in this difference. The corresponding reconstruction residual is
\begin{equation}
    \vec{r}_{\mathrm{err},i}^{m}(x,c;\theta,\phi)
    =\Delta\vec{u}_{i}^{m}(x,c;\theta)
     -\widehat{\Delta\vec{u}}_{K,i}^{\ell\rightarrow m}(x,c;\theta,\phi).
    \label{eq:reconstruction-residual}
\end{equation}
The subscript $\mathrm{err}$ identifies a reconstruction residual, not a second causal mechanism. Because the full change, selected reconstruction, and residual are inserted in separate nonlinear forward passes, their answer effects are not assumed to add.

\subsection{Shared residual-state subspaces}

Sparse reconstructions and native channels use particular fitted or axis-aligned coordinates. To test whether the clean-minus-masked hidden-state change is distributed across coordinates but lies in a low-dimensional space shared across examples within one model, we fit a separate orthonormal basis for each model. For a frozen layer $\ell$ and position group, let $\vec{B}_{r}^{\ell}\in\mathbb{R}^{d\times r}$ contain the first $r$ orthonormal directions obtained by uncentered principal-component analysis of the clean-minus-masked differences from a designated fitting split. Section~\ref{sec:experimental-setup} fixes the fitting, selection, and replication roles. For fixed $(x,c;\theta)$, the projection of Equation~\ref{eq:hidden-difference} onto this basis is
\begin{equation}
    \Delta\vec{h}_{r,i}^{\ell}
    =\vec{B}_{r}^{\ell}(\vec{B}_{r}^{\ell})^{\mathsf{T}}
      \Delta\vec{h}_{i}^{\ell}.
    \label{eq:shared-subspace-projection}
\end{equation}
The subscript $r$ is the number of shared basis directions, and $i$ ranges over the frozen position group. Restoration adds $\Delta\vec{h}_{r,i}^{\ell}$ to the masked run, whereas deletion subtracts it from the clean run. The full-hidden reference uses the unprojected $\Delta\vec{h}_{i}^{\ell}$ at the same layer and positions. The fitted basis is an analysis object constructed from intervention differences; it is not a native variable used by the VLM during generation.

\subsection{Common intervention design}

The primary visual treatment is corruption of the human-annotated \emph{answer mask}, which covers the localized region supporting the short answer. When answer-bearing evidence and necessary adjacent context are separately annotated, their union provides a broader evidence condition. The primary control is a same-area random mask on the same image. Gray masking is the common corruption operator; shifted, shuffled, blurred, inpainted, patch-shuffled, dilated, and eroded variants are auxiliary checks reported in the appendix. Ambiguous or diffuse evidence is excluded rather than forced into a compact mask.

Internal interventions evaluate five distinct object classes. \emph{Hidden-state interventions} restore or delete complete residual-stream states at predefined layers and token-position groups. \emph{Sparse-component interventions} insert or remove a selected sparse reconstruction or its reconstruction residual at one matched feed-forward output. \emph{Native-channel interventions} act on coordinates after the gated feed-forward calculation and before its output projection. \emph{Native-channel path tests} jointly manipulate upstream and downstream channel groups. \emph{Shared-subspace interventions} restore or delete the projected hidden-state difference in Equation~\ref{eq:shared-subspace-projection}. This ordering separates an internally recoverable effect, the part visible to fitted sparse coordinates, compact groups of native coordinates, and a compact non-axis-aligned alternative before testing the complete hidden state.

Each result subsection specifies the locations and internal objects needed for its question. Hidden-state locations are compared through a predefined position grid in Section~\ref{sec:shared-dependence}. Sparse components are tested at a matched feed-forward output in Section~\ref{sec:sparse-effect}; Section~\ref{sec:sparse-structure} then describes their normalized structures without treating them as faithful paths. Section~\ref{sec:native-path} tests native-channel candidates without fitted-tool reconstruction error. Section~\ref{sec:subspace-bottleneck} evaluates shared subspaces and the frozen LLaVA full-hidden candidate, ending the intervention sequence with the independent full-hidden confirmation and matched follow-up.

The next section fixes the model--interface pairs, evaluation sets, target construction, selection roles, and statistical criteria before any result is interpreted.
